% TODO make hw stylesheet
\documentclass[10pt]{article}
\usepackage[letterpaper,top=1in,left=1in,right=1in]{geometry}
\usepackage[dvipsnames,svgnames]{xcolor}
\usepackage[shortlabels]{enumitem}
\usepackage[framemethod=TikZ]{mdframed}
\usepackage{amsmath,amssymb,amsthm}
\usepackage[colorlinks]{hyperref}
\usepackage{multicol}
\usepackage{tabularx}
\usepackage{bbm}
\usepackage{tikz-cd}

\hypersetup{
    colorlinks=true,
    linkcolor=blue,
    filecolor=magenta,
    urlcolor=magenta,
    pdftitle={MAT 535 HW}
}

\usepackage{mathtools}
\usepackage{thmtools}
\usepackage[textsize=footnotesize]{todonotes}
\usepackage{titling}
\usepackage{cancel}

\reversemarginpar
\mdfdefinestyle{mdbluebox}{roundcorner=10pt,innerbottommargin=9pt,
linecolor=blue,backgroundcolor=TealBlue!5,}
\declaretheoremstyle[headfont=\sffamily\bfseries\color{MidnightBlue},
mdframed={style=mdbluebox},]{thmbluebox}

\mdfdefinestyle{mdbloobox}{frametitlefont=\bfseries,innerbottommargin=8pt,
nobreak=true,backgroundcolor=TealBlue!5,linecolor=blue,}
\declaretheoremstyle[headfont=\bfseries\color{MidnightBlue},
mdframed={style=mdbloobox},headpunct={\\[3pt]},postheadspace=0pt,]{thmbloobox}

\mdfdefinestyle{mdredbox}{frametitlefont=\bfseries,innerbottommargin=8pt,
nobreak=true,backgroundcolor=Salmon!5,linecolor=RawSienna,}
\declaretheoremstyle[headfont=\bfseries\color{RawSienna},
mdframed={style=mdredbox},headpunct={\\[3pt]},postheadspace=0pt,]{thmredbox}

\mdfdefinestyle{mdgreenbox}{linecolor=ForestGreen,backgroundcolor=ForestGreen!5,
linewidth=2pt,rightline=false,leftline=true,topline=false,bottomline=false,}
\declaretheoremstyle[headfont=\bfseries\sffamily\color{ForestGreen!70!black},
mdframed={style=mdgreenbox},headpunct={ --- },]{thmgreenbox}

\mdfdefinestyle{mdorangebox}{linecolor=RedOrange,backgroundcolor=RedOrange!5,
linewidth=2pt,rightline=false,leftline=true,topline=false,bottomline=false,}
\declaretheoremstyle[headfont=\bfseries\sffamily\color{RedOrange!70!black},
mdframed={style=mdorangebox},headpunct={ --- },]{thmorangebox}

\mdfdefinestyle{mdblackbox}{linecolor=black,backgroundcolor=RedViolet!5!gray!5,
linewidth=3pt,rightline=false,leftline=true,topline=false,bottomline=false,}
\declaretheoremstyle[mdframed={style=mdblackbox}]{thmblackbox}

\declaretheoremstyle[spaceabove=6pt,spacebelow=6pt]{boldhead}
\declaretheoremstyle[spaceabove=6pt,spacebelow=6pt,headfont=\normalfont\itshape]{italichead}

\declaretheorem[style=boldhead,name=Problem]{problem}
\declaretheorem[style=boldhead,name=Problem,numbered=no]{problem*}
\declaretheorem[style=boldhead,name=Setup]{setup} % for config geo
\declaretheorem[style=boldhead,name=Example,sibling=problem]{example}
\declaretheorem[style=boldhead,name=$\bigstar$ Problem,sibling=problem]{niceprob}
\declaretheorem[style=boldhead,name=Theorem,sibling=problem]{theorem}
\declaretheorem[style=boldhead,name=Corollary,sibling=theorem]{corollary}
\declaretheorem[style=boldhead,name=Proposition,sibling=theorem]{prop}
\declaretheorem[style=boldhead,name=Lemma,numberwithin=problem]{lemma}
\declaretheorem[style=boldhead,name=Theorem,numbered=no]{theorem*}
\declaretheorem[style=boldhead,name=Lemma,numbered=no]{lemma*}
\declaretheorem[style=boldhead,name=Claim,numberwithin=problem]{claim}
\declaretheorem[style=boldhead,name=Claim,numbered=no]{claim*}
\declaretheorem[style=boldhead,name=Remark,numberwithin=problem]{remark}
\declaretheorem[style=boldhead,name=Remark,numbered=no]{remark*}
\declaretheorem[style=boldhead,name=Definition,sibling=theorem]{definition}
\declaretheorem[style=boldhead,name=Definition,numbered=no]{definition*}

\providecommand{\ol}{\overline}
\renewcommand{\hom}{\operatorname{Hom}}
\providecommand{\ceil}[1]{\left\lceil #1 \right\rceil}
\providecommand{\floor}[1]{\left\lfloor #1 \right\rfloor}
\newcommand{\dd}{\,\mathrm{d}}
\providecommand{\ul}{\underline}
\providecommand{\eps}{\varepsilon}
\providecommand{\half}{\frac{1}{2}}
\providecommand{\CC}{\mathbb C}
\providecommand{\FF}{\mathbb F}
\providecommand{\II}{\mathbb I}
\providecommand{\NN}{\mathbb N}
\providecommand{\QQ}{\mathbb Q}
\providecommand{\RR}{\mathbb R}
\providecommand{\ZZ}{\mathbb Z}
\providecommand{\ind}{\mathbbm 1}
\providecommand{\dg}{^\circ}
\providecommand{\ii}{\item}
\providecommand{\setdiff}{\smallsetminus}
\providecommand{\alert}[1]{{\sffamily\textbf{\textcolor{blue}{#1}}}}
\providecommand{\inv}{^{-1}}
\providecommand{\mb}[1]{\mathbf{#1}}
\newcommand{\what}{\widehat}

\newenvironment{soln}{\begin{proof}[Solution]}{\end{proof}}

\providecommand{\pow}{\operatorname{Pow}}
\providecommand{\vocab}[1]{{\textbf{\textcolor{ForestGreen}{#1}}}}
\providecommand{\lcm}{\operatorname{lcm}}
\providecommand{\abs}[1]{\left\lvert #1\right\rvert}
\providecommand{\smabs}[1]{\lvert #1\rvert}
\providecommand{\innerprod}[1]{\langle\!\langle #1\rangle\!\rangle}
\providecommand{\norm}[1]{\left\| #1\right\|}
\providecommand{\smnorm}[1]{\| #1\|}
\providecommand{\gr}{\operatorname{gr}}

\usepackage{mathrsfs}

% place title at top right
\setlength{\droptitle}{-8ex}
\pretitle{\begin{flushright}\Large\bfseries}
\posttitle{\par\end{flushright}}
\preauthor{\begin{flushright}}
\postauthor{\end{flushright}}
\predate{\begin{flushright}}
\postdate{\end{flushright}}

\title{MAT 535 HW02}
\author{\large Aditya Pahuja}
\date{\large Due 02/12}
\begin{document}
\maketitle

\begin{problem}
    [\S10.4\#14]
    Let $I$ be a nonempty index set and, for each $i\in I$,
    let $N_i$ be a left $R$-module. Let $M$ be a right $R$-module.
    Prove the group isomorphism $M\otimes(\bigoplus_{i\in I} N_i)
    \cong\bigoplus_{i\in I}(M\otimes N_i)$.
\end{problem}

\begin{soln}
    Define the map $\phi\colon M\times(\bigoplus_{i\in I} N_i)
    \to \bigoplus_{i\in I}(M\otimes N_i)$ by
    \begin{equation*}
	\phi(m,(n_i)_{i\in I}) = (m\otimes n_i)_{i\in I}.
    \end{equation*}
    This map is clearly additive in each argument and
    \[
	\phi(mr,(n_i)_{\in I}) = (mr\otimes n_i)_{i\in I}
	= (m\otimes rn_i)_{i\in I} = \phi(m,r(n_i)_{i\in I}),
    \]
    so $\phi$ is $R$-balanced, thus inducing a homomorphism
    $\widetilde\phi\colon M\otimes(\bigoplus_{i\in I}N_i)\to
    \bigoplus_{i\in I}(M\otimes N_i)$ which is uniquely determined by
    $\widetilde\phi(m\otimes (n_i)_{i\in I}) = (m\otimes n_i)_{i\in I}$.

    In the other direction, for each $\alpha\in I$, the map
    $\psi_\alpha\colon M\times N_\alpha\to M\otimes(\bigoplus_{i\in I}N_i)$
    given by $\psi_\alpha(m,n) = m\otimes\iota_\alpha(n)$,
    where $\iota_\alpha\colon N_\alpha\to\bigoplus_{i\in I}N_i$
    is the canonical embedding of $N_\alpha$ into the direct sum,
    is $R$-balanced. This means $\psi_\alpha$ induces
    a homomorphism $\widetilde\psi_\alpha\colon M\otimes N_\alpha
    \to M\otimes(\bigoplus_{i\in I}N_i)$ which is determined by
    $m\otimes n\mapsto m\otimes\iota_\alpha(n)$.
    Then, the $\widetilde\psi_\alpha$ canonically induce a map
    $\widetilde\psi\colon\bigoplus_{i\in I}(M\otimes N_i)\to
    M\otimes(\bigoplus_{i\in I}N_i)$ determined by
    \begin{equation*}
	\widetilde\psi((m_i\otimes n_i)_{i\in I})
	= \sum_{i\in I}m_i\otimes \iota_i(n_i)
    \end{equation*}
    where $n_i\in N_i$ and cofinitely many of the
    $m_i\otimes n_i$, hence cofinitely many of
    the $m_i\otimes\iota_i(n_i)$ summands, are zero.

    We show that $\widetilde\phi$ and $\widetilde\psi$ are inverses.
    For sequences of simple tensors
    $(m_i\otimes n_i)_{i\in I}$ in $\bigoplus_{i\in I}(M\otimes N_i)$
    with $m_i\otimes n_i = 0$ for cofinitely many $i\in I$,
    \begin{equation*}
        \widetilde\phi\circ
	\widetilde\psi((m_i\otimes n_i)_{i\in I})
	= \sum_{i\in I}\widetilde\phi(m_i\otimes\iota_i(n_i))
	= \sum_{i\in I} m_i\otimes n_i
	= (m_i\otimes n_i)_{i\in I}
    \end{equation*}
    (To be precise, the summands in the latter sum are really
    elements of the embedding of $M\otimes N_i$ into the direct sum
    of tensors.) A generic element of the direct sum of $M\otimes N_i$
    is given by a finite sum of simple tensors at finitely many indices
    $i\in I$, which can be written as the sum of elements
    of the form $(m_i\otimes n_i)_{i\in I}$ as in the equation above,
    so $\widetilde\phi\circ\widetilde\psi$ is the identity on
    $\bigoplus_{i\in I}(M\otimes N_i)$. Conversely,
    for simple tensors $m\otimes(n_i)_{i\in I}\in M\otimes
    (\bigoplus_{i\in I} N_i)$,
    \begin{equation*}
	\widetilde\psi\circ\widetilde\phi(m\otimes(n_i)_{i\in I})
	= \widetilde\psi((m\otimes n_i)_{i\in I})
	= \sum_{i\in I}m\otimes\iota_i(n_i)
	= m\otimes\left(\sum_{i\in I}\iota_i(n_i)\right)
	= m\otimes(n_i)_{i\in I},
    \end{equation*}
    implying that $\widetilde\psi\circ\widetilde\phi$
    is the identity on $M\otimes(\bigoplus_{i\in I}N_i)$,
    and proving that $\widetilde\phi$ and $\widetilde\psi$
    are isomorphisms.
\end{soln}

\pagebreak
\begin{problem}
    [\S10.4\#16]
    Suppose $R$ is commutative and let $I$ and $J$ be ideals of $R$,
    so $R/I$ and $R/J$ are naturally $R$-modules.
    \begin{enumerate}[(a)]
        \ii Prove that every element of $R/I\otimes_R R/J$
	can be written as a simple tensor of the form
	$(1 + I)\otimes(r + J)$.
	\ii Prove that there is an $R$-module isomorphism
	$R/I\otimes_R R/J\cong R/(I+J)$
	mapping $(r+I)\otimes(r'+J)$ to $rr'+(I+J)$.
    \end{enumerate}
\end{problem}

\begin{soln}
    (a) Any simple tensor can be written as
    $(a + I)\otimes(b + J) = (1 + I)\otimes(ab + J)$,
    so a generic element of $R/I\otimes_R R/J$,
    which is a sum of finitely many simple tensors, can be written as
    \begin{equation*}
	\sum_{i=1}^m (1 + I)\otimes(r_i + J)
	= (1 + I)\otimes\left(\sum_{i=1}^m r_i + J\right).
    \end{equation*}

    (b) It is easy to see that the map
    $(r+I,r'+J)\mapsto rr'+(I+J)$ from $R/I\times R/J$
    to $R/(I+J)$ is $R$-balanced (since $(sr+I,r'+J)$ and
    $(r+I,sr'+J)$ both map to $srr' + (I+J)$), so the given map,
    call it $\phi$, is a well-defined homomorphism of $R$-modules.

    We will show that the map $\psi\colon R/(I+J)\to R/I\otimes_RR/J$
    defined by $\psi(r+(I+J)) = (1+I)\otimes(r+J)$
    is the inverse of $\phi$, so that $\phi$ is indeed
    an isomorphism of $R$-modules. First, to see that $\psi$ is
    well-defined, we note that the map
    $R\to R/I\otimes_RR/J$ defined by $r\mapsto(1+I)\otimes(r+J)$
    has both $I$ and $J$, hence $I+J$, as a subset of its kernel, since
    $r\in I$ implies $(1+I)\otimes(r+J) = (r+I)\otimes(1+J)
    = 0_{R/I}\otimes(1+J) = 0$ and $r\in J$ implies $r+J = 0_{R/J}$,
    so this map on $R$ factors through the quotient $R/(I+J)$
    to yield $\psi$. Now we check that $\phi$ and $\psi$ are inverses:
    \begin{gather*}
        \psi\circ\phi((1+I)\otimes(r+J))
	= \psi(r + (I+J)) = (1 + I)\otimes(r+J),\\
	\phi\circ\psi(r+(I+J)) = \phi((1 + I)\otimes(r+J))
	= r + (I+J)
    \end{gather*}
    so $\psi\circ\phi$ and $\phi\circ\psi$ are respectively the
    identities on $R/I\otimes_RR/J$ and $R/(I+J)$, where we use part (a)
    to ensure that we only need to check that simple tensors of the form
    $(1+I)\otimes(r+J)$ get sent to themselves by $\psi\circ\phi$.
    Thus, $\phi$ is indeed an isomorphism of the given $R$-modules.
\end{soln}

\begin{problem}
    [\S10.4\#21]
    Suppose  $R$ is commutative and let $I$ and $J$ be ideals of $R$.
    \begin{enumerate}[(a)]
        \ii Show that there is a surjective $R$-module homomorphism
	from $I\otimes_R J$ to the product ideal $IJ$
	mapping $i\otimes j$ to the element $ij$.
	\ii Give an example to show that the map in (a)
	need not be injective.
    \end{enumerate}
\end{problem}

\begin{soln}
    (a) The map $I\times J\to IJ$ given by $(i,j)\mapsto ij$
    is clearly additive in both arguments
    and satisfies $(ir,j) = (i,rj)$, so it is $R$-balanced;
    thus, the map $\phi\colon I\otimes_RJ\to IJ$ given by
    $\phi(i\otimes j) = ij$ is a well-defined homomorphism of $R$-modules.
    Recall that $IJ$ is generated as an $R$-module
    by finite sums of the form
    $\sum_{k=1}^m i_kj_k$ for $i_k\in I$, $j_k\in J$ for $k=1,\dots,m$,
    so it suffices to show that all such finite sums
    are in the image of $\phi$.
    Indeed, for a given sum $\sum_{k=1}^m i_kj_k$,
    \begin{equation*}
	\phi\left(\sum_{k=1}^m i_k\otimes j_k\right)
	= \sum_{k=1}^m i_kj_k,
    \end{equation*}
    so $\phi(I\otimes_RJ) = IJ$.
    \\

    (b) Exercise 17 in this section shows that,
    with $R = \ZZ[x]$ and $I = J = (2,x)$,
    the elements $2\otimes x,x\otimes2\in I\otimes_RI$ are distinct,
    while $\phi(2\otimes x) = \phi(x\otimes 2) = 2x$.
    Thus, $\phi$ has nontrivial kernel, for example
    containing the nonzero element $2\otimes x - x\otimes 2$.
\end{soln}

\pagebreak
\begin{problem}
    [\S10.5\#2]
    Suppose that
    \begin{center}
	\begin{tikzcd}
	    A\arrow[r,"f"]\arrow[d,"\alpha"]&
	    B\arrow[r,"g"]\arrow[d,"\beta"] &
	    C\arrow[r,"h"]\arrow[d,"\gamma"] &
	    D\arrow[d,"\delta"]\\
	    A'\arrow[r,"f'"]&
	    B'\arrow[r,"g'"]&
	    C'\arrow[r,"h'"]&
	    D'
	\end{tikzcd}
    \end{center}
    is a commutative diagram of groups and that the rows are exact.
    Prove that
    \begin{enumerate}[(a)]
        \ii if $\alpha$ is surjective, and $\beta$, $\delta$ are injective,
	then $\gamma$ is injective.
	\ii if $\delta$ is injective, and $\alpha$, $\gamma$ are surjective,
	then $\beta$ is surjective.
    \end{enumerate}
\end{problem}

\begin{soln}
    (a) Suppose $c\in\ker\gamma$; we will show $c=1$.
    Then $\delta h(c) = 1$ and $\delta$ is injective,
    so $h(c) = 1$, meaning $c\in\ker h = \operatorname{im}g$.
    Write $c = g(b)$ for $b\in B$; then
    \begin{equation*}
        1 = \gamma(c) = \gamma g(b) = g'\beta(b),
    \end{equation*}
    so $\beta(b) \in\ker g' = \operatorname{im}f'$.
    Write $\beta(b) = f'(a')$ for $a'\in A'$;
    since $\alpha$ is surjective, $a' = \alpha(a)$ for some $a\in A$.
    Thus $\beta(b) = f'\alpha(a) = \beta f(a)$;
    by the injectivity of $\beta$, this means $f(a) = b$,
    i.e. $b\in\operatorname{im}f = \ker g$.
    Therefore $c = g(b) = 1$, i.e. $\ker\gamma$ consists of
    only the identity.\\

    (b) Suppose $b'\in B'$. By surjectivity,
    $g'(b') = \gamma(c)$ for $c\in C$. By exactness at $C'$,
    $\gamma(c) = g'(b')\in\ker h'$, hence
    \begin{equation*}
        1 = h'\gamma(c) = \delta h(c),
    \end{equation*}
    implying by injectivity that $h(c) = 1$.
    By exactness at $C$, $c = g(b)$ for some $b\in B$. Observe that
    \begin{equation*}
        g'(\beta(b)\cdot(b')\inv)
	= \gamma g(b)\cdot g'(b')\inv
	= \gamma(c)\cdot \gamma(c)\inv = 1,
    \end{equation*}
    so $\beta(b)\cdot (b')\inv\in\ker g' = \operatorname{im}f'$,
    hence equals $f'(a')$ for some $a'\in A'$;
    surjectivity of $\alpha$ then implies that $a' = \alpha(a)$
    for some $a\in A$. That is,
    $\beta(b)\cdot(b')\inv = f'\alpha(a) = \beta f(a)$
    so $\beta(b\cdot f(a)\inv) = b'$,
    meaning $b'$ is in the image of $\beta$.
\end{soln}

\begin{problem}
    [\S10.5\#5]
    Prove that an arbitrary direct sum $\bigoplus_{i\in I}A_i$
    of $R$-modules is flat if and only if each $A_i$ is flat.
\end{problem}

\begin{soln}
    To show flatness of an $R$-module $A$, it suffices to show that
    for an injective $R$-module homomorphism $f\colon L\to M$,
    $f\otimes 1\colon L\otimes_R A\to M\otimes_RA$ is also injective,
    so fix an injective map $f\colon L\to M$.

    If $\bigoplus_{i\in I}A_i$ is flat, for a given $R$-module $B$
    and index $i\in I$ let
    $\iota_{B,i}\colon B\otimes A_i\to B\otimes_R(\bigoplus_{i\in I}A_i)$
    be the canonical embedding of $B\otimes_R A_i$
    into $\bigoplus_{i\in I}(B\otimes_RA_i)$,
    composed with the isomorphism
    $\phi_B\colon \bigoplus_{i\in I}(B\otimes_RA_i)\to
    B\otimes_R(\bigoplus_{i\in I}A_i)$ defined in the first problem.
    Note that $\iota_{B,i}$ is injective for each $B$
    and that we have a commutative diagram
    \begin{center}
        \begin{tikzcd}
	    L\otimes_RA_i\arrow[r,"f\otimes 1"]\arrow[d,"\iota_{L,i}"] &
	    M\otimes_RA_i\arrow[d,"\iota_{M,i}"]\\
	    L\otimes_R(\bigoplus_{i\in I}A_i)\arrow[r,"f\otimes 1"]&
	    M\otimes_R(\bigoplus_{i\in I}A_i).
        \end{tikzcd}
    \end{center}
    Since the morphism $(f\otimes 1)\circ\iota_{L,i}$
    is equal to $\iota_{M,i}\circ(f\otimes 1)$
    and the former, being a composition of injective morphisms,
    is injective, the latter morphism must also be injective,
    implying that the upper $f\otimes 1$ is injective,
    showing that $-\otimes A_i$ is left-exact for each $i\in I$.

    Conversely, if $A_i$ is flat for each $i\in I$,
    then $f\otimes 1_i\colon L\otimes_RA_i\to M\otimes_RA_i$
    is injective for each $i$ where $1_i$ is the identity on $A_i$.
    The injective maps induce injective maps between the coproducts,
    which, upon passing through $\phi_L$ and $\phi_N$,
    shows that $f\otimes 1\colon L\otimes_R(\bigoplus_{i\in I}A_i)\to
    M\otimes_R(\bigoplus_{i\in I}A_i)$ is injective.
\end{soln}

\begin{problem}
    This exercises proves Theorem 38 that
    every left $R$-module is contained in an injective left $R$-module.
    \begin{enumerate}[(a)]
        \ii Show that $M$ is contained in an injective $\ZZ$-module $Q$.
	\ii Show that
	$\hom_R(R,M)\subseteq\hom_\ZZ(R,M)\subseteq\hom_\ZZ(R,Q)$.
	\ii Use the $R$-module isomorphism $M\cong\hom_R(R,M)$
	and the previous exercise to conclude that
	$M$ is contained in an injective $R$-module.
    \end{enumerate}
\end{problem}

\begin{soln}
    (a) This follows immediately from Corollary 37,
    which says that the problem holds when $R = \ZZ$,
    because all left $R$-modules are also $\ZZ$-modules.
    \\

    (b) The first inclusion holds because
    an $R$-module homomorphism $\phi\colon R\to M$
    is also a homomorphism of abelian groups
    if we just ignore the multiplicative structure of the map.
    
    The second inclusion holds because we have an injective map
    sending morphisms $f\in\hom(R,M)$ to $\iota\circ f\in\hom(R,Q)$ where
    $\iota\colon M\hookrightarrow Q$ is the inclusion of $M$ into $Q$.
    \\

    (c) The preceding exercise shows that
    $\hom_\ZZ(R,Q)$ is an injective $R$-module if $Q$ is,
    so the chain of inclusions
    $R\cong\hom_R(R,M)\hookrightarrow\hom_\ZZ(R,M)
    \hookrightarrow\hom_\ZZ(R,Q)$ from part (b)
    implies that $R$ is isomorphic to a submodule
    of the injective left $R$-module $\hom_\ZZ(R,Q)$
    as desired.
\end{soln}










\end{document}
